\documentclass
  [hyperref={colorlinks = true,linkcolor = blue, 
             citecolor = blue, urlcolor = blue}
  ]{beamer}

\setbeamertemplate{navigation symbols}{}
\setbeamertemplate{footline}[frame number]
\setbeamercolor{itemize item}{fg=white}
 
\usepackage[utf8]{inputenc}

\usepackage[newfloat]{minted}

\usepackage{pgf}
\usepackage{tikz}
\usepackage{upquote}
\usepackage{natbib}
\usepackage[export]{adjustbox}
\usepackage{amsmath}
\usepackage{amssymb}
\usepackage{trfrac}

\usetikzlibrary{arrows,automata,fit, shapes.geometric}

\bibliographystyle{abbrvnat}

\newenvironment{code}{\captionsetup{type=listing}}{}
\SetupFloatingEnvironment{listing}{name=Listing}

\setbeamertemplate{items}[square]

\title{Introduction to Propositions as Types}
\author{Donovan Crichton}
  
\date{March 2024}

\begin{document}
 
\frame{\titlepage}

\begin{frame}[fragile]
  \frametitle{Preliminaries}
  \begin{itemize}
  \item Slides and Examples available at:
    \href{https://github.com/donovancrichton/Talks}
         {https://github.com/donovancrichton/Talks}
  \item This talk: BFPG/PropositionsAsTypes
  \end{itemize}
\end{frame}

\begin{frame}[fragile]
\frametitle{About me}
\begin{minipage}[t]{\linewidth}
\noindent
  \begin{minipage}{0.5\linewidth}
    \begin{figure}[h]
       \includegraphics[width=\linewidth]
         {Images/ANU.png}
    \end{figure}
  \end{minipage}
  \hfill
  \begin{minipage}{0.45\linewidth}
    \begin{itemize}
      \item 
        \href{https://comp.anu.edu.au/people/donovan-crichton/}
             {PhD Candidate}
      \item 
        \href{https://comp.anu.edu.au/research/clusters/computing-foundations/}
             {Computing Foundations}
      \item 
        \href{https://comp.anu.edu.au/}
             {School of Computing}
    \end{itemize}
  \end{minipage}
\end{minipage}
\hrule
  \begin{minipage}{0.5\linewidth}
    \begin{figure}[h]
      \includegraphics[scale=0.5,width=\linewidth]
        {Images/GriffithLogo.png}
    \end{figure}
  \end{minipage}
  \hfill
  \begin{minipage}{0.45\linewidth}
  \begin{itemize}
    \item Visiting Scholar
    \item Trusted Systems Lab
    \item \href{https://www.griffith.edu.au/institute-integrated-intelligent-systems}{IIIS}
  \end{itemize}
  \end{minipage}
\hrule
  \begin{minipage}{0.5\linewidth}
    \begin{figure}[h]
      \includegraphics[width=\linewidth]
        {Images/asd-logo.png}
    \end{figure}
  \end{minipage}
  \hfill
  \begin{minipage}{0.45\linewidth}
  \begin{itemize}
    \item ASD 
      \href{https://www.asd.gov.au/about/asd-anu-co-lab}{Co-Lab} 
      Scholar
  \end{itemize}
  \end{minipage}
\end{frame}

% Okay, what do I want to cover here?
% I want to introduce the concept of a proposition
% I want to introduce the concept of a theorem
% I want to introduce the concept of a type
% I want to introduce the curry-howard isomorphism
% I want to talk about isomorphsim vs equality
% 

\begin{frame}[fragile]
  \frametitle{Proving Theorems $\cong$ Writing Programs - Overview}
  \setbeamercolor{block title}{fg=white,bg=green!55!black}
  \setbeamercolor{block body}{fg=white,bg=green!25!black}
  \begin{block}{The PAT (Propositions as Types) Interpretation}
    \begin{itemize}
      \item Types are logical propositions or theorems.
      \item An inhabitant or element of the type is the proof.
      \item Strictly speaking an isomorphism that holds under specific
      conditions.
    \end{itemize}
  \end{block}
  % Why is this important?
  \setbeamercolor{block title}{fg=white,bg=orange!55!black}
  \setbeamercolor{block body}{fg=white,bg=orange!25!black}
  \begin{block}{Uses for the PAT Interpretation}
    \begin{itemize}
      \item Used by mathematicians and logicians to \textit{formalise} or
            \textit{mechanise} mathematics and logic.
      \item Used by software developers to provide a stronger guarantee of
            functional correctness than testing.
    \end{itemize}
  \end{block}
\end{frame}

\begin{frame}[fragile]
  \frametitle{Proving Theorems $\cong$ Writing Programs - Applications}
  % Who uses this in anger?
  \setbeamercolor{block title}{fg=white,bg=orange!55!black}
  \setbeamercolor{block body}{fg=white,bg=orange!25!black}
  \begin{block}{Applications}
    \begin{itemize}
      \item In defense/security: guarantees for software security.
      \item In aeronautics and naval agencies: guarantees for software
      properties that control hardware.
      \item In fin-tech: guarnatees for software properties that involve
      transactions.
      \item Suitable for any application where the cost of a failed test
            after deployment is just too high.
    \end{itemize}
  \end{block}

\end{frame}


\begin{frame}[fragile]
  \frametitle{Propostions - Refresher}
% What is a proposition
  \setbeamercolor{block title}{fg=white,bg=green!55!black}
  \setbeamercolor{block body}{fg=white,bg=green!25!black}
  \begin{block}{What is a proposition?}
  A statement or assertion that expresses a judgement, usually readily
  apparent or verifiable.
  \end{block}
  \setbeamercolor{block title}{fg=white,bg=orange!55!black}
  \setbeamercolor{block body}{fg=white,bg=orange!25!black}
  \begin{block}{Examples of a proposition}
    \begin{itemize}
      \item ``It is raining outside.''
      \item ``Terry is the child of Leigh.''
      \item Sometimes formally denoted: 
            \\ Let $p$ denote ``It is raining outside.''
    \end{itemize}
  \end{block}
\end{frame}

\begin{frame}[fragile]
  \frametitle{Theorems - Refresher}
% What is a theorem
  \setbeamercolor{block title}{fg=white,bg=green!55!black}
  \setbeamercolor{block body}{fg=white,bg=green!25!black}
  \begin{block}{What is a theorem?}
  A more general proposition that is less readily-apparent, usually requires
  a chain of reasoning to be accepted, or 'proven'.
  \end{block}
  \setbeamercolor{block title}{fg=white,bg=orange!55!black}
  \setbeamercolor{block body}{fg=white,bg=orange!25!black}
  \begin{block}{Examples of Theorems}
    \begin{itemize}
      \item ``If Terry has a child $r$, and Terry is the child of Leigh, then
      $r$ is the grandchild of Leigh.''
      \item ``For all natural numbers $x$ and $y$, and given the addition
      operation ($+$), then $x + y$ is equal to $y + x$.''
      \item ``There is at least one weekday occuring in the future.''
    \end{itemize}
  \end{block}
\end{frame}

\begin{frame}[fragile]
  \frametitle{Proofs - Refresher}
% What is a proof
  \setbeamercolor{block title}{fg=white,bg=green!55!black}
  \setbeamercolor{block body}{fg=white,bg=green!25!black}
  \begin{block}{What is a poof?}
    The chain of formal reasoning that, when followed, always verifies a
    theorem.
  \end{block}
  \setbeamercolor{block title}{fg=white,bg=orange!55!black}
  \setbeamercolor{block body}{fg=white,bg=orange!25!black}
  \begin{block}{Examples of Proof sketches}
    \begin{itemize}
      \item Grandchild is defined as "The child of a person's child." Thus,
      under this definition, this holds.
      \item Use a lemma (a smaller proof) to show that $x + 0 = 0 + x$ by
      definition of $+$. Then apply the induction hypothesis for all other
      cases.
      \item This is more complicated. We could use an inductive, discrete definition of
      ``future'', and show that after so many iterations a new day occurs.
    \end{itemize}
  \end{block}
\end{frame}



\begin{frame}[fragile]
  \frametitle{Types - Refresher}
% What is a type?
  \setbeamercolor{block title}{fg=white,bg=green!55!black}
  \setbeamercolor{block body}{fg=white,bg=green!25!black}
  \begin{block}{What is a type?}
  We often think of types in two forms, one computational, and one
  mathematical.
    \begin{itemize}
      \item A type is a way to interpret the series of bits that represent a
      value to distinguish it from other values that are encoded with the same
      bit representation. 
      \item A type is a ``set'' of elements that inhabit the type.
    \end{itemize}
  \end{block}
  \setbeamercolor{block title}{fg=white,bg=orange!55!black}
  \setbeamercolor{block body}{fg=white,bg=orange!25!black}
  \begin{block}{Examples of Types}
    \begin{itemize}
      \item Take the binary number 1100001. Does this represent the positive
      integer 1,100,001? The letter `a' on the keyboard? Or the decimal value
      of 97?
      \item $\mathbb{B}$ denotes the set of Boolean values: \{True, False\}.
      \item $\mathbb{Z}$ denotes the set of Integer values: \{..., -2, -1, 0,
      1, 2, ...\}.
    \end{itemize}
  \end{block}
\end{frame}

\begin{frame}[fragile]
  \frametitle{Elements/Inhabitants - Refresher}
% What is an inhabitant of a type
  \setbeamercolor{block title}{fg=white,bg=green!55!black}
  \setbeamercolor{block body}{fg=white,bg=green!25!black}
  \begin{block}{What is an element or inhabitant of a type?}
    One of the members of the set that characterises a type.
  \end{block}
  \setbeamercolor{block title}{fg=white,bg=orange!55!black}
  \setbeamercolor{block body}{fg=white,bg=orange!25!black}
  \begin{block}{Examples of elements of types.}
    \begin{itemize}
      \item Ordinary values: Such as \{True,False\} from $\mathbb{B}$.
      \item Data Types:  \\ 
        \{(``Nil'', $\ast \rightarrow$ List A), (``Cons'', A
      $\rightarrow$ List A $\rightarrow$ List A)\}.
      \item Codata Types: \{(``hd'', List A $\rightarrow$ A), (``tl'', List A
      $\rightarrow$ List A)\}.
    \end{itemize}
  \end{block}
\end{frame}

\begin{frame}[fragile]
  \frametitle{The CHL Correspondance - Computational Trinitarianism}
  % draw a graph with arrows between Curry, Howard, and Lambeck
\end{frame}

\begin{frame}[fragile]
  \frametitle{Isomorphism vs Equality}
  % define isomorphsim and equality
\end{frame}

\begin{frame}[fragile]
  \frametitle{Intuitionistic Logic $\cong$ STLC}
  % Point out that it can hold for languages isomorphic to STLC 
\end{frame}

\begin{frame}[fragile]
  \frametitle{Types and Props - Implication}
\end{frame}

\begin{frame}[fragile]
  \frametitle{Types and Props - Conjuction}
\end{frame}

\begin{frame}[fragile]
  \frametitle{Types and Props - Disjunction}
\end{frame}

\begin{frame}[fragile]
  \frametitle{Types and Props - Negation}
\end{frame}

\begin{frame}[fragile]
  \frametitle{Types and Props - Quantifiers?}
\end{frame}

\begin{frame}[fragile]
  \frametitle{Types and Props - Quantifiers (and Language)}
\end{frame}

\begin{frame}[fragile]
  \frametitle{Examples - Linked List Invariants and Vectors}
\end{frame}

\begin{frame}[fragile]
\frametitle{References}
\bibliography{references}{}
\end{frame}

\end{document}

